\subsection*{A 1P1C matmul skeleton}

Event-based synchronization is left out of this skeleton; Section~\ref{sec:3-8-5} adds the full pipeline protocol. Focus on the program structure:

\begin{lstlisting}[language=croqtile,caption={A 1P1C matmul skeleton}]
__co__ void matmul(global f16 [M, K] lhs, global f16 [N, K] rhs, global f16 [M, N] output) {
  parallel {block_m, block_n} by [cdiv(M, WARP_M), cdiv(N, WARP_N)] : block {
    shared f16 [WARP_M, TILE_K] lhs_load_s;
    shared f16 [WARP_N, TILE_K] rhs_load_s;
    shared f16 [WARP_M, WARP_N] output_s;
    parallel p1 by 2 : group-4 {
      inthreads.async (p1 == 0) {
        foreach {iv_k} in [cdiv(K, TILE_K)] {
          dma.copy lhs.subspan(WARP_M, TILE_K).at(block_m, iv_k) => lhs_load_s;
          dma.copy rhs.chunkat(block_n, iv_k) => rhs_load_s;
        }
      }
      inthreads.async (p1 == 1) {
        mc = mma.fill.f16 0.0f;
        foreach {iv_k} in [cdiv(K, TILE_K)] {
          foreach {iv_warp} in [cdiv(TILE_K, WARP_K)] {
            ma = mma.load lhs_load_s.chunkat(_, iv_warp);
            mb = mma.load rhs_load_s.chunkat(_, iv_warp);
            mma.row.row mc, ma, mb;
          }
        }
        mma.store mc, output_s;
        dma.copy output_s => output.subspan(WARP_M, WARP_N).at(block_m, block_n);
      }
    }
  }
}
\end{lstlisting}

\textbf{Producer \texttt{foreach}} --- Walks K with \texttt{cdiv(K, TILE\_\allowbreak{}K)} steps;

\textbf{Consumer \texttt{mma} path}

\textbf{Missing coordination} --- Both sides loop over K independently. The consumer assumes each K-slab is ready when it reads. Making that assumption correct requires events (Section~\ref{sec:3-8-5}).
